\chapter{Background Review}
\label{ch:background}

This chapter follows a four-level stack of design decisions that ascend from material to control. Table~\ref{tab:ch2-overview} summarizes the structure: each section addresses one level of the body and one corresponding question that the literature on that level has tried to answer.

\begin{table}[h]
\centering
\footnotesize
\caption{Four-level structure of Chapter~\ref{ch:background}.}
\label{tab:ch2-overview}
\begin{tabular}{@{}l l p{0.45\linewidth}@{}}
\toprule
\textbf{Section} & \textbf{Level} & \textbf{Field and core question} \\
\midrule
\S\ref{sec:soft-robotics}   & Material        & Soft robotics and pneumatic actuation: how does a deformable body change what actuation can do? \\
\addlinespace
\S\ref{sec:shape-changing}  & Structural      & Morphological computation: how does shape turn deformation into purposeful motion? \\
\addlinespace
\S\ref{sec:modularity}      & Configurational & Modular reconfigurable robotics: how does the arrangement of comparable units become a design variable? \\
\addlinespace
\S\ref{sec:emergence}       & Control         & Emergence and bottom-up coordination: how can complex behavior arise from simple modules? \\
\bottomrule
\end{tabular}
\end{table}

Within each section I begin with a conceptual breakdown of the level it occupies, then read a handful of representative robotic systems that instantiate that concept, and close by reflecting on what these systems offer this thesis and what they leave undone. Across the four sections, those small reflections accumulate into a single research site and a single set of gaps to which the rest of the thesis responds.
